% sec-01: the release.  Figure 13 (mermaid-type gantt), Table 21.
\section{The Release}

Six checks precede any transmission, and any one of them failing stops the entry.
They are restated here rather than referenced, because a check nobody can find is
a check nobody runs.

The market is open, confirmed against the exchange calendar rather than assumed.
The release approval was given on the preceding day and is unchanged overnight.
The financing branch is unchanged, so the queue stands as cut against the
nine-month horizon rather than being amended after entry. The auction calendar
has been re-read this morning, because auctions do not settle on a federal
holiday and a calendar read before one is wrong for the week after it
\cite{treasuryauctions2026}. The limit prices were written this morning from the
last completed session's close. And the reserve balance is current, since every
order size is a share of it.

\begin{appfloat}[!tb]
\begin{appfig}[mermaid-type / gantt across one market session]
% A seven-hour rule at 1.55cm per hour, nine bars on a 0.52cm pitch.  \ganttrow
% draws one bar, so a bar is moved by changing one number.  The auction cutoff
% is a vertical rule rather than a bar, because a cutoff is an instant.
\node[mmband,minimum width=0.39cm,minimum height=5.4cm,anchor=south west] at (0,-4.95) {};
\node[mmband,minimum width=0.39cm,minimum height=5.4cm,anchor=south west] at (10.46,-4.95) {};
\axisx{0}{10.85}{0.42}
\foreach \x/\lab in {0/{open}, 1.55/{+1h}, 3.10/{+2h}, 4.65/{+3h}, 6.20/{+4h},
                     7.75/{+5h}, 9.30/{+6h}, 10.85/{close}}{
  \draw[fundgray,line width=0.3pt] (\x,0.42) -- (\x,0.28);
  \node[font=\tiny\sffamily,text=fundgray,anchor=south] at (\x,0.48) {\lab};}
\ganttrow{0}{0}{0.65}{mmbarg}{Confirm open, re-read auctions}
\ganttrow{-0.52}{0.65}{1.04}{mmbarg}{Write this morning's limits}
\ganttrow{-1.04}{1.30}{1.82}{mmbark}{Release four held letters}
\ganttrow{-1.56}{1.82}{2.34}{mmbark}{Enter auction bids, lines 1 to 3}
\ganttrow{-2.08}{2.45}{2.71}{mmbar}{Enter line 4, limit, day}
\ganttrow{-2.60}{3.23}{3.49}{mmbar}{Enter line 5, limit, day}
\ganttrow{-3.12}{3.62}{4.40}{mmbar}{Send the three follow-ups}
\ganttrow{-3.64}{4.65}{5.81}{mmbar}{Submit the portal items}
\ganttrow{-4.16}{9.30}{9.82}{mmbark}{Reconcile, write the fill record}
\ganttrow{-4.68}{9.82}{10.21}{mmbark}{File the release record}
\draw[fundaccent,line width=0.9pt] (2.34,0.30) -- (2.34,-4.95);
\node[font=\tiny\sffamily,text=fundaccent,anchor=west] at (2.42,0.62) {auction cutoff, an instant};
\legkey{0}{-5.45}{fundgrayl}{Before the open}
\legkey{2.60}{-5.45}{fundmid}{Release and record}
\legkey{5.60}{-5.45}{fundpale}{Entry and submission}
\pnote{7.90}{-5.45}{Shaded bands: the first and last fifteen minutes.}
\end{appfig}
\vspace{-0.60cm}
\figcaption{Figure 13. One open market session, with the release and execution steps in\\
sequence and the auction cutoff drawn as an instant that constrains them all.}
\end{appfloat}

\subsection{What was released, and what was withheld}

\begin{apptable}
\begin{tabularx}{\textwidth}{>{\raggedright\arraybackslash}p{4.2cm} >{\raggedright\arraybackslash}p{2.1cm} >{\raggedright\arraybackslash}X}
\headrow \thead{Item} & \thead{Released} & \thead{The dependency, or the reason it was withheld}\\
\midrule
Congressional district letter & Yes & Held from the closed day; district offices are open \\
State finance center letter & Yes & The same \\
Regional development letter & Yes & The same \\
FDA Pre-Request for Designation & Yes & The same. The most consequential item in the block \\
Federal follow-up, round two & Yes & Three business days elapsed since the day 1 letters \\
Developer follow-up, window only & Yes & Three business days elapsed; one question, not two \\
Six queued orders & Yes & Six checks passed, in the sequence an auction cutoff sets \\
Portal submissions & Yes & Support desks open again \\
Site follow-up & \textbf{No} & Only two business days have elapsed since the day 3 letters \\
Foundation follow-up & \textbf{No} & The same \\
Site meeting confirmation & Conditional & Sent only if an institution replied; otherwise left unsent \\
\end{tabularx}
\end{apptable}
\vspace{-0.60cm}
\tabcap{Table 21. Eleven items with whether each was released, and the interval\\
or dependency that decided it, including the three deliberately withheld.}

\subsection{Why the withheld items matter more than the released ones}

A block that releases everything it has written is a block with no discipline in
it. The federal holiday is not a business day, which makes this the third
business day after day 1 and only the second after day 3. A site that receives a
follow-up on the second business day reads a company counting hours, and the
impression outlasts the letter.

The conditional item is the same principle in a different form. A confirmation
letter with nothing to confirm is a second request wearing a confirmation's
clothes, and recipients notice.
